\chapter{Related Work}
\label{chap:related_work}

Chapter~\ref{chap:introduction} framed this work around the cross-attack
robustness gap: adversarial robustness can depend strongly on the attack used
during training, and a single robustness score can hide uneven behavior across
threat models. This chapter reviews the related work through that lens. The
goal is not to list every adversarial training variant, but to identify which
parts of the literature motivate the method developed in
Chapter~\ref{chap:methodology} and which assumptions remain insufficient for
the practical multi-attack setting studied in this report.

The chapter proceeds in five steps. Section~3.1 reviews how adversarial
training moved from single-threat training toward multi-attack and
multi-perturbation objectives. Section~3.2 discusses Group DRO as a mechanism
for emphasizing weak groups, while also highlighting its dependence on useful
group definitions. Section~3.3 connects the attacks-as-domains view to domain
generalization, but distinguishes standard domain generalization from
threat-model-aware adversarial evaluation. Section~3.4 turns to hidden-group
and cluster-discovery methods, which motivate latent grouping but also expose
the limits of representation-only clusters. Section~3.5 synthesizes these
threads into the gap addressed by AttackDRO++.

\section{Multi-Attack Adversarial Training}

% TODO[Thành]:
% Goal:
% - Explain how prior work trains against multiple attacks and where uniform or fixed attack mixing falls short.
% Flow:
% - Summarize the main strategy families, then compare them to this report's Multi-AT baseline.
% - End with the limitation that attack diversity alone does not discover hard latent groups.
% Include:
% - Representative works only where they support the gap argument.
% Avoid:
% - Paper-by-paper listing without synthesis.
% Target length: 4-6 paragraphs.

%Table comparing: attacks used, objective function, sampling strategy
%Tramer 2019, Maini 2020 (MSD), Madaan 2021 (DPP), Laidlaw 2021




\section{Group DRO Applied to Robustness}

%Sagawa 2020 (seminal)
%Liu 2021, Zhai 2022 extensions
%Connection to fairness (Hashimoto 2018)


\section{Domain Generalization Meets Adversarial Robustness}



\section{Cluster Discovery in Robust Training Pipelines}

% TODO[Thành]:
% Goal:
% - Review clustering or representation-based group discovery methods relevant to robust training.
% Flow:
% - Explain why learned representations can reveal hidden groups.
% - Compare representation-only grouping with the gradient-fingerprint idea.
% - End with the limitation that prior grouping may not capture optimization-level difficulty.
% Include:
% - Mention K-means only as needed for the report's method.
% Avoid:
% - A broad clustering survey.
% Target length: 3-5 paragraphs.

\section{The Remaining Gap and Motivation for This Work}

% TODO[Thành]:
% Goal:
% - Synthesize the previous related-work sections into the exact gap this report addresses.
% Flow:
% - Multi-attack training gives diversity but not adaptive hard-group emphasis.
% - Group DRO gives adaptive emphasis but depends on useful group definitions.
% - Domain/clustering ideas suggest group discovery, but robustness needs attack-aware evaluation.
% Include:
% - A clear final paragraph that makes Chapter 4 feel necessary.
% Avoid:
% - Introducing new papers that were not discussed earlier.
% Target length: 3-4 paragraphs.

